% Generated draft: requires mathematical and source review.
\documentclass[11pt]{article}
\usepackage[margin=25mm]{geometry}
\usepackage{fontspec}
\setmainfont{Latin Modern Roman}
\usepackage{amsmath,amssymb,mathtools,unicode-math}
\setmathfont{Latin Modern Math}
\usepackage{enumitem,needspace,xurl,hyperref}
\setlength{\parindent}{0pt}
\setlength{\parskip}{5pt}
\setlength{\emergencystretch}{3em}
\newcommand{\N}{\mathbb{N}}
\newcommand{\Z}{\mathbb{Z}}
\newcommand{\Q}{\mathbb{Q}}
\newcommand{\R}{\mathbb{R}}
\newcommand{\C}{\mathbb{C}}
\newcommand{\F}{\mathbb{F}}
\newcommand{\A}{\mathbb{A}}
\newcommand{\PP}{\mathbb{P}}
\newcommand{\E}{\mathbb{E}}
\newcommand{\eps}{\varepsilon}
\newcommand{\dd}{\,\mathrm{d}}
\begin{document}
\section*{UnsolvedMath ID: 30005096}
\textbf{Source problem code:} catalogue entry
\subsection{Definitions and mathematical statement}
Let \(n\geq 1\) be an integer and let \(\mathbb H=\{\tau\in\C:\operatorname{Im}\tau>0\}\). The modular curve \(Y(1)\) has complex analytic realization \(\operatorname{SL}_2(\Z)\backslash\mathbb H\); the modular \(j\)-invariant identifies its underlying algebraic curve with \(\A^1_{\Q}\). For the formulation below, use its complex base change \(Y(1)_{\C}\). This is an explicit interpretation of the analytic setting: the supplied clean statement does not specify a field of definition for \(V\).

A singular modulus is a number \(j(\tau)\) with \(\tau\in\mathbb H\) imaginary quadratic over \(\Q\). For \(N\geq1\), let \(\Phi_N(X,Y)\) be the classical modular polynomial, whose vanishing expresses the existence of a cyclic isogeny of degree \(N\) between elliptic curves with the indicated \(j\)-invariants. A special subvariety of \(Y(1)_{\C}^{n}\) is an irreducible component of a locus defined by finitely many equations of the forms
\[
 x_i=c,\qquad \Phi_N(x_i,x_j)=0,
\]
where \(c\) is a singular modulus and \(N\) is a positive integer. Here special subvarieties, rather than arbitrary weakly special subvarieties, are intended. An irreducible algebraic subvariety \(V\subseteq Y(1)_{\C}^{n}\) is Hodge generic if it is contained in no proper special subvariety.

For a proper special subvariety \(S\subsetneq Y(1)_{\C}^{n}\), an irreducible component \(W\) of the algebraic intersection \(V\cap S\) is called atypical if
\[
 \dim W>\dim V+\dim S-n.
\]
Define
\[
 \mathcal A(V)=
 \bigcup_{\substack{S\subsetneq Y(1)_{\C}^{n}\\S\text{ special}}}
 \ \bigcup_{\substack{W\text{ an irreducible component of }V\cap S\\
 \dim W>\dim V+\dim S-n}}W.
\]
The supplied question is whether, for every \(n\geq1\) and every irreducible Hodge-generic algebraic subvariety \(V\subseteq Y(1)_{\C}^{n}\),
\[
 \overline{\mathcal A(V)}^{\,\mathrm{Zar},V}\subsetneq V.
\]
Equivalently, must there exist a proper Zariski-closed subset of \(V\) containing every such atypical component?

The general target above comes from the supplied clean statement. The supplied original excerpt is truncated and discusses only an irreducible Hodge-generic curve \(C\) defined over \(\Q\), with an additional asymmetric hypothesis for a reported theorem. It neither defines asymmetric nor supplies a complete statement of that theorem. Consequently, the excerpt alone does not establish the intended field-of-definition scope of the general question or permit a precise formulation of the reported special case.
\subsection{Short English statement}
Must all unexpectedly large components of intersections between an irreducible Hodge-generic subvariety of a power of the modular curve and proper special subvarieties lie in one proper Zariski-closed subset? The displayed formulation uses complex varieties explicitly; the source leaves the intended field-of-definition restriction unresolved.
\subsection{Sources}
\begin{itemize}
\item Supplied local UnsolvedMath registry record, permanent ID \(30005096\), title: Zilber--Pink for Powers of the Modular Curve; exact source code: \texttt{catalogue entry}. The registry labels the problem \texttt{partially\_solved} and its statement \texttt{corrected\_verified}. Its assertion that the statement was checked against the official report is supplied provenance, not independent verification in this draft.
\item Supplied report citation: \emph{Diophantische Approximationen} (2023), Oberwolfach Reports, pp. 1091--1164, DOI 10.4171/owr/2022/21. \url{https://doi.org/10.4171/owr/2022/21}. The record gives proposed year 2022. The excerpt attributes the contribution \emph{Zilber--Pink in \(Y(1)^n\)} to Christopher Daw, joint work with Martin Orr; no separate proposer is supplied.
\item Supplied literature source: Daw, Zilber--Pink in \(Y(1)\) powers. \url{https://doi.org/10.4171/owr/2022/21}.
\item Supplied literature source: Daw--Orr, beyond multiplicative degeneration. \url{https://arxiv.org/abs/2402.09487}.
\item Literature assessment dated exactly \texttt{2026-08-22}: the supplied curation states that Zilber--Pink for asymmetric curves is a theorem, that later work adds boundary cases, and that the general problem in powers of the modular curve remains open. This is supplied dated curation derived from the OWR material, not an independently verified assessment. The background expressly requests independent verification of time-sensitive status and cited claims.
\item The original excerpt reports a 2012 result of Habegger and Pila for irreducible Hodge-generic asymmetric curves defined over \(\Q\), citing its reference [5]. The supplied record does not contain the bibliographic entry for [5] or the definition of asymmetric.
\end{itemize}
\end{document}
